Physical Informed Neural Networks (PiNNs) are a recent approach that exploit physical laws inside the learning process of a Neural Network.
As previously stated, we have used a bottom-up approach by first testing this approach tu a Unicycle. This robot still has a nonlinear dynamics but it's way simpler than the quadrotor one. Its configuration vector is composed by the cartesian position $(x,y)$ and its orientation angle $\theta$. 
\begin{align}
    \begin{cases}
        \dot{x} = v \cos(\theta) \\
        \dot{y} = v \sin(\theta) \\
        \dot{\theta} = w \\
    \end{cases}
\end{align}

\begin{align}
    \mathbf{R}_{\text{RPY}} &=
        \begin{bmatrix}
        \cos\psi \cos\theta & \cos\psi \sin\theta \sin\phi - \sin\psi \cos\phi & \cos\psi \sin\theta \cos\phi + \sin\psi \sin\phi \\
        \sin\psi \cos\theta & \sin\psi \sin\theta \sin\phi + \cos\psi \cos\phi & \sin\psi \sin\theta \cos\phi - \cos\psi \sin\phi \\
        -\sin\theta & \cos\theta \sin\phi & \cos\theta \cos\phi
        \end{bmatrix}
    \\
    \dot{\mathbf{v}} &= 
        \begin{bmatrix}
    0.0 \\
    0.0 \\
    -\text{g}
        \end{bmatrix}
    +
    \mathbf{R}_{\text{RPY}} \cdot
    \begin{bmatrix}
        0.0 \\
        0.0 \\
        \frac{T}{m}
    \end{bmatrix}
\end{align}
In our implementation, we have used a simplified dynamic model for the quadrotor instead of the simplified second order system.
\begin{align}
    \begin{cases}
        \ddot{x} = -(\cos\psi \sin\theta \cos\phi + \sin\psi \sin\phi) \frac{T}{m}\\
        \ddot{y} = -(\sin\psi \sin\theta \cos\phi - \cos\psi \sin\phi) \frac{T}{m} \\
        \ddot{z} = -(\cos\theta \cos\phi) \frac{T}{m} +g\\
        \ddot{\phi} = \frac{\tau_{\phi}}{I_{x}}\\
        \ddot{\theta} = \frac{\tau_{\theta}}{I_{y}}\\
        \ddot{\psi} = \frac{\tau_{\psi}}{I_{z}}\\ 
    \end{cases}  
\end{align}

\begin{align}
    \begin{cases}
        \dot{x} = v_{x}\\
        \dot{y} = v_{y}\\
        \dot{z} = v_{z}\\
        \dot{v}_{x} = (\cos\psi \sin\theta \cos\phi + \sin\psi \sin\phi) \frac{T}{m}\\
        \dot{v}_{y} = (\sin\psi \sin\theta \cos\phi - \cos\psi \sin\phi) \frac{T}{m} \\
        \dot{v}_{z} = (\cos\theta \cos\phi) \frac{T}{m} -g\\
        \dot{\phi} = \\
        \dot{\theta} = \\
        \dot{\psi} = \\ 
    \end{cases}  
\end{align}
In our implementation, we have used a PiNNs to predict the next state of the system, given the control input obtained from the Model Predictive Control.
\todo[inline]{losses}
\todo[inline]{pipeline}
\todo[inline]{image pinn}